\section{Conclusion}
\label{sec:conclusion}
This thesis focuses on access control for a hospital management system (HMS) built with a microservices architecture.
In a hospital, data must be protected because it contains sensitive patient information. At the same time, staff need fast
access to perform medical tasks. This creates a practical tension between security and availability.
\vspace{0.5cm}


Many hospital systems use role-based access control (RBAC) because it is easy to understand and manage. Roles such as
doctor, nurse, pharmacist, and receptionist match how hospitals are organized. However, RBAC alone is often too rigid.
In real workflows, access depends on context: department assignment, care team relationship, shift time, or whether a staff
member is responsible for a specific patient. Attribute-based access control (ABAC) can describe these situations, but a
pure ABAC approach can become difficult to maintain when policies grow.
\vspace{0.5cm}


To handle this, this thesis proposes a hybrid RBAC/ABAC model following a role-centric direction (RBAC-A). In this design,
roles define the baseline permissions, and attributes are used mainly as constraints to narrow access in specific contexts.
This keeps policies readable while still supporting common hospital conditions.
\vspace{0.5cm}


The thesis also shows how the model can be applied in a microservices HMS without spreading authorization logic across
many services. The implementation uses a centralized enforcement pattern. The API gateway acts as the Policy Enforcement Point
(PEP) and blocks or forwards requests. A dedicated Authorization Service acts as the Policy Decision Point (PDP) and
evaluates requests using Open Policy Agent (OPA). Policies are managed through a Policy Service, and audit logs are recorded
to support traceability.
\vspace{0.5cm}


From the implementation and validation work, this thesis provides three main takeaways. First, role-centric hybrid policies
are easier to review than fully attribute-centric policies, because the baseline permission is explicit in each rule.
Second, separating policy evaluation from business services reduces duplicated checks and makes enforcement consistent.
Third, audit logging is not an optional feature in a hospital setting; it is necessary for accountability and later review.
\vspace{0.5cm}


Overall, this thesis demonstrates that a role-centric hybrid RBAC/ABAC approach is a practical middle ground for hospital
microservices. It keeps the structure of RBAC while adding enough context checks to fit real hospital workflows.
\vspace{0.5cm}


\section{Limitations}
\label{sec:limitations}
The demo in this thesis is designed to be clear and testable, but it still has limitations.
\vspace{0.5cm}


First, the implementation covers a limited set of services and workflows. A full HMS in production typically includes more
modules and more data flows, which can introduce additional policy cases. Second, the policy set and test scenarios represent
typical assumptions, but real hospitals may have special procedures and local compliance requirements. Third, the conflict
detection approach in the application is practical and helpful during policy updates, but it is not a full formal verification
method. When conditions become more complex, conflict reasoning can also become more complex.
\vspace{0.5cm}


These limitations suggest that the current system should be seen as a foundation. It shows the feasibility of the model and
the architecture, but more work is needed for large-scale, real-world deployment.

\section{Future Work}
\label{sec:future_work}
Future work can improve the system in several directions. The most important point is to extend the model carefully while
keeping policies understandable and easy to manage.
\vspace{0.5cm}


A first direction is to support more hospital scenarios with richer context. For example, policies can include stronger
relationship checks (doctor--patient assignment and care team membership), location constraints (on-site vs remote access),
and purpose-of-use constraints (treatment vs billing). In addition, data sensitivity levels can be used to apply different
rules or different disclosure levels, such as masking some fields for non-primary staff.
\vspace{0.5cm}


A second direction is emergency access (break-glass). In hospitals, there are situations where strict rules must be
temporarily bypassed to protect a patient. A safe break-glass design should not simply ``allow everything''. It should
require a clear trigger, collect a justification, record a stronger audit trail, and support post-event review. This makes
emergency access controlled and accountable.
\vspace{0.5cm}


A third direction is to improve the policy lifecycle and tooling. In practice, policies change over time and are often
managed by multiple administrators. Future work can add policy versioning, review/approval steps, and a staging environment
before publishing to production. It is also useful to build automated policy tests, where expected Permit/Deny scenarios are
stored and re-run after every policy update. For conflict detection, the system can provide better reports and suggestions,
such as narrowing a condition or adjusting priority when a conflict is found.
\vspace{0.5cm}


A fourth direction is performance and scalability. As the number of rules and requests grows, decision latency becomes more
important. Future work can study caching strategies, reduce repeated attribute lookups, and optimize how policy bundles are
distributed to OPA instances. Load testing at larger scales should also be performed with realistic mixes of endpoints.
\vspace{0.5cm}


Finally, future work can strengthen assurance and compliance alignment. This includes clearer policy combining behavior,
more systematic testing, and better reporting for compliance review. If the system is integrated with healthcare standards
(e.g., resource modeling aligned with common healthcare formats), it can reduce integration cost in real environments.

\section{Final Remarks}
\label{sec:final_remarks}
This thesis proposed a role-centric hybrid RBAC/ABAC model for hospital microservices and demonstrated it through a working
demonstration with centralized enforcement at the gateway and centralized decision evaluation using OPA. The approach keeps the
simplicity of RBAC for baseline permissions and uses attributes as constraints to match dynamic hospital context. With further
work on emergency access, policy tooling, and large-scale evaluation, this design can be extended toward a deployable access
control solution for hospital systems.